% Punkte und Strecken im Koordinatensystem · Vierecke nachweisen · Prüfungs-Fokus Abitur GK – bank/punkte-und-strecken-im-koordinatensystem/e4.jsonl (zusammenbau v0.6, Prüfungs-Fokus)
\einheitenkopf[e4]{Einheit 4 · Vierecke nachweisen}
\zweigzeile{Abi GK ×9}
\setcounter{aufgabe}{0}

% Anlauf (ohne Prüfkennung)
\begin{aufgabe}{Vierecke nachweisen – Anlauf}
\begin{teile}
\teil Nachzuweisen: $ABCD$ ist ein Rechteck. Was gehört zum Steckbrief? \\ \kreuz{vier gleich lange Seiten} \\ \kreuz{Parallelogramm und ein rechter Winkel} \\ \kreuz{ein Paar paralleler Seiten}
\end{teile}
\end{aufgabe}

\begin{aufgabe}{Vierecke nachweisen – Prüfungsaufgaben}
\begin{teile}
\teil Zeige: $A(2 | 1 | 0)$, $B(4 | 3 | 1)$, $C(2 | 4 | 3)$, $D(0 | 2 | 2)$ bilden ein Rechteck $ABCD$. \hfill (A23)
\teil $A(1 | 5 | p)$, $B(4 | 5 | p)$, $C(4 | 1 | 5)$ und $D(1 | 1 | 5)$ mit einer natürlichen Zahl $p$. Zeige: $ABCD$ ist für jedes $p$ ein Rechteck. Länge von $BC$ in Abhängigkeit von $p$? \hfill (A25)
\teil Zeige: $A(0 | 1 | 3)$, $B(2 | 3 | 4)$, $C(3 | 5 | 6)$, $D(1 | 3 | 5)$ bilden eine Raute $ABCD$. \hfill (A22)
\teil Eine Dachfläche hat die Ecken $A(0 | 1 | 1)$, $B(4 | 1 | 4)$, $C(4 | 6 | 4)$, $D(0 | 6 | 1)$. Zeige: $ABCD$ ist eine Raute. \hfill (A22)
\teil Zeige: $A(1 | 2 | 2)$, $B(2 | 6 | 10)$, $C(6 | 7 | 18)$, $D(5 | 3 | 10)$ bilden eine Raute $ABCD$. \hfill (A22)
\teil Zeige: $A(2 | 2 | 2)$, $B(5 | 3 | 2)$, $C(6 | 5 | 4)$, $D(3 | 4 | 4)$ bilden ein Parallelogramm, aber kein Rechteck. \hfill (A22)
\teil Zeige: $A(1 | 3 | 2)$, $B(3 | 4 | 4)$, $C(5 | 2 | 5)$, $D(3 | 1 | 3)$ bilden eine Raute, aber kein Quadrat. \hfill (A20)
\teil Zeige: $A(1 | 0 | 1)$, $B(1 | 3 | 5)$, $C(4 | 7 | 5)$, $D(4 | 4 | 1)$ bilden eine Raute, aber kein Quadrat. \hfill (A20)
\teil Eine Kletterwand hat die Ecken $A(4 | 1 | 3)$, $B(1 | 4 | 3)$, $F(1 | 7 | 0)$ und $E(7 | 1 | 0)$ (1 LE = 1 m). Zeige: $ABFE$ ist ein Trapez, in dem zwei gegenüberliegende Seiten gleich lang sind. \hfill (A18)
\teil Ein Würfel mit Kante $8$ hat eine Ecke im Ursprung und liegt im ersten Oktanten. $P(8 | 0 | 2)$, $Q(3 | 8 | 0)$, $R(0 | 8 | 3)$ und $S(2 | 0 | 8)$ liegen auf seinen Kanten. Begründe: $PQRS$ ist ein Trapez mit zwei gleich langen Seiten, und $PS$ ist doppelt so lang wie $QR$. \hfill (A19)
\end{teile}
\end{aufgabe}

\begin{aufgabe}{Vierecke nachweisen – Prüfungsaufgaben – weiter}
\begin{teile}
\teil Eine Terrasse ist das Viereck $A(1 | -3 | 0)$, $B(7 | -1 | 0)$, $C(7 | 1 | 0)$, $D(1 | 3 | 0)$ (1 LE = 1 m). Zeige: $ABCD$ ist ein Trapez. Flächeninhalt? \hfill (A26) \\ \leerfeld[m²]
\teil Ein Beet ist das Viereck $E(-2 | 1 | 0)$, $F(6 | 1 | 0)$, $G(3 | 5 | 0)$, $H(1 | 5 | 0)$ (1 LE = 1 m). Zeige: $EFGH$ ist ein Trapez. Flächeninhalt? \hfill (A26) \\ \leerfeld[m²]
\teil Eine Glasplatte ist das Viereck $K(0 | -3 | 1)$, $L(4 | -1 | 1)$, $M(4 | 1 | 1)$, $N(0 | 3 | 1)$ (1 LE = 1 dm). Zeige: $KLMN$ ist ein Trapez. Flächeninhalt? \hfill (A26) \\ \leerfeld dm²
\teil Zeige: Das Viereck $K(0 | 0 | 2)$, $L(2 | 3 | 2)$, $M(6 | 0 | 2)$, $N(2 | -3 | 2)$ ist ein Drachenviereck. \hfill (A23)
\teil Zeige: Das Viereck $K(1 | 2 | 0)$, $L(1 | 5 | 2)$, $M(1 | 2 | 8)$, $N(1 | -1 | 2)$ ist ein Drachenviereck. \hfill (A23)
\teil Zeige: Das Viereck $P(6 | 1 | 1)$, $Q(6 | 5 | 1)$, $R(6 | 4 | 4)$, $S(6 | 1 | 5)$ ist ein Drachenviereck. \hfill (A23)
\teil Eine ebene Rasenfläche hat die Ecken $A(2 | 1 | 0)$, $B(9 | 1 | 0)$, $C(8 | 4 | 0{,}5)$, $D(8 | 7 | 1)$ und $E(2 | 7 | 1)$ (1 LE = 1 m). Zeige: $AE$ ist parallel zu $CD$, und $CD$ und $DE$ bilden einen rechten Winkel. \hfill (A21)
\teil Zeige: Im Viereck $A(1 | -1 | 4)$, $B(3 | 1 | 5)$, $C(1 | 2 | 2{,}5)$, $D(0 | 1 | 2)$ ist $AB$ parallel zu $DC$, und bei $A$ liegt ein rechter Winkel. \hfill (A21)
\teil Zeige: Im Viereck $A(0 | 2 | 1)$, $B(6 | 2 | 9)$, $C(7 | 3 | 2)$, $D(4 | 3 | -2)$ ist $AB$ parallel zu $DC$, und bei $A$ liegt ein rechter Winkel. \hfill (A21)
\teil $P(2 | 1 | 0)$ und $Q(2 | 1 | 5)$ sind benachbarte Ecken eines Quadrats in der Ebene $3x_1 + 4x_2 = 10$. Gib eine weitere Ecke an. \hfill (A22)
\teil $P(3 | 1 | 4)$ und $Q(3 | 6 | 4)$ sind benachbarte Ecken eines Quadrats in der Ebene $4x_1 - 3x_3 = 0$. Gib eine weitere Ecke an. \hfill (A22)
\end{teile}
\end{aufgabe}

\begin{aufgabe}{Vierecke nachweisen – Prüfungsaufgaben – weiter}
\begin{teile}
\teil Das Viereck $ABCD$ hat die Ecken $A(0 | 0 | 0)$, $B(3 | 1 | 0)$, $C(0 | 7 | 0)$ und $D(-2 | 3 | 0)$. $B$ wird parallel zur $x_2$-Achse nach $B'$ verschoben, sodass $AB'CD$ bei $B'$ einen rechten Innenwinkel hat; $b$ ist die gesuchte Koordinate. Erläutere den Ansatz $(3 | b | 0) \circ (3 | b - 7 | 0) = 0$. \hfill (A25)
\teil Das Viereck $ABCD$ hat die Ecken $A(0 | 0 | 0)$, $B(2 | 4 | 0)$, $C(9 | 0 | 0)$ und $D(4 | -3 | 0)$. $B$ wird parallel zur $x_1$-Achse nach $B'$ verschoben, sodass $AB'CD$ bei $B'$ einen rechten Innenwinkel hat; $b$ ist die gesuchte Koordinate. Erläutere den Ansatz $(b | 4 | 0) \circ (b - 9 | 4 | 0) = 0$. \hfill (A25)
\teil Die Raute $A(5 | 3 | 1)$, $B(2 | 3 | 5)$, $C(-1 | 3 | 1)$, $D(2 | 3 | -3)$ hat einen Inkreis, der alle vier Seiten berührt. Begründe: $P(3{,}92 | 3 | 2{,}44)$ ist einer der Berührpunkte. \hfill (A20)
\teil Die Raute $A(1 | 1 | 6)$, $B(4 | 1 | 2)$, $C(1 | 1 | -2)$, $D(-2 | 1 | 2)$ hat einen Inkreis, der alle vier Seiten berührt. Begründe: $P(2{,}92 | 1 | 3{,}44)$ ist einer der Berührpunkte. \hfill (A20)
\teil Eine Wand hat die Form des Parallelogramms $A(2 | 1 | 0)$, $B(9 | 1 | 0)$, $C(10 | 1 | 4)$, $D(3 | 1 | 4)$ (1 LE = 1 m). Ein Farbstreifen von $A$ nach $P(6 | 1 | 4)$ auf der Seite $DC$ teilt sie. Welcher Teil ist größer? Begründe ohne Flächenberechnung. \hfill (A22)
\teil Eine rechteckige Wand hat die Ecken $A(1 | 3 | 0)$, $B(1 | 9 | 0)$, $C(1 | 9 | 5)$, $D(1 | 3 | 5)$ (1 LE = 1 m). Ein Streifen von $B$ nach $Q(1 | 3 | 2)$ auf der Seite $AD$ teilt sie. Welcher Teil ist größer? Begründe ohne Flächenberechnung. \hfill (A22)
\end{teile}
\end{aufgabe}

\medskip
\uebersichtskasten{Steckbriefe (Ecken ABCD der Reihe nach um das Viereck): Parallelogramm: AB = DC. Rechteck: Parallelogramm und AB ∘ AD = 0. Raute: Parallelogramm und alle vier Seiten gleich lang. Quadrat: Raute mit rechtem Winkel. Trapez: AB = k · DC (kollineare Gegenseiten) – der Faktor k liefert das Längenverhältnis der parallelen Seiten. Drachenviereck: zwei Paare benachbarter gleich langer Seiten. \\ A(0 | 0 | 0), B(4 | 2 | 0), C(5 | 6 | 1), D(1 | 4 | 1): AB = (4 | 2 | 0) = DC → Parallelogramm; AB ∘ AD = 4 + 8 + 0 = 12 ≠ 0 → kein Rechteck. \\ Reihenfolge: Gegenseiten sind AB und DC, nicht CD – ein Vorzeichenfehler macht aus dem Parallelogramm ein unmögliches Viereck. \\ Vollständig nachweisen, knapp ausschließen: im Raum sichern vier gleiche Seitenlängen allein die Raute nicht – erst das Parallelogramm macht das Viereck eben, dann tragen die Längen; zum Ausschließen genügt dagegen ein einziges Gegenmerkmal (Skalarprodukt ungleich null → kein Rechteck, kein Quadrat). \\ Flächeninhalt als Anhang: Trapez A = ½ · (a + c) · h, Dreieck A = ½ · g · h – die Höhe aus den Koordinaten (achsenparallele Lage) oder als Abstand paralleler Seiten.}
